\section{Second-Order Expansion of the BCA Conditions}

\subsection{Production and Resource Constraint}

For the Cobb-Douglas technology in Equation \eqref{eq:production},
\begin{equation}
    \log Y_t
    =
    a_t+\alpha\log K_t+(1-\alpha)\log L_t
    \label{eq:log_production_exact}
\end{equation}
is exact. Therefore, TFP level shocks enter the efficiency wedge directly. TFP
uncertainty does not create a separate production wedge at date \(t\); it changes
the conditional distribution of future \(A_{t+1}\).

The resource constraint is nonlinear in log coordinates. If
\(\hat{y}_t=\log(Y_t/\bar{Y})\), \(\hat{c}_t=\log(C_t/\bar{C})\), and
\(\hat{x}_t=\log(X_t/\bar{X})\), then
\begin{equation}
    \exp(\hat{y}_t)
    =
    s_c\exp(\hat{c}_t)
    +
    s_x\exp(\hat{x}_t)
    +
    s_g\exp(\hat{g}_t),
    \label{eq:exact_resource_log_coords}
\end{equation}
where \(s_j=\bar{J}/\bar{Y}\) for \(J\in\{C,X,G\}\). A second-order expansion
gives
\begin{equation}
    \hat{y}_t+\frac{1}{2}\hat{y}_t^2
    =
    s_c\left(\hat{c}_t+\frac{1}{2}\hat{c}_t^2\right)
    +
    s_x\left(\hat{x}_t+\frac{1}{2}\hat{x}_t^2\right)
    +
    s_g\left(\hat{g}_t+\frac{1}{2}\hat{g}_t^2\right)
    +
    O(3).
    \label{eq:second_order_resource}
\end{equation}
These curvature terms are resource-accounting terms. They should not be
confused with price-setting or demand wedges.

\subsection{Labor Wedge}

The labor wedge is static:
\[
    -\frac{U_L(C_t,L_t)}{U_C(C_t,L_t)}
    =
    (1-\tau_{\ell,t})A_tF_L(K_t,L_t).
\]
With the separable utility in Equation \eqref{eq:period_utility_example}, the
log labor wedge is
\begin{equation}
    \log(1-\tau_{\ell,t})
    =
    \log\chi+\nu\log L_t+\gamma\log C_t
    -
    \log A_t
    -
    \log F_L(K_t,L_t).
    \label{eq:log_labor_wedge}
\end{equation}
Because this condition is static, there is no direct Jensen term involving
\(\Var_t(a_{t+1})\). TFP uncertainty can move the measured labor wedge only
indirectly through equilibrium movements in \(C_t\), \(L_t\), \(K_t\), and
\(A_t\), or through model misspecification when the accounting system is too
small.

\subsection{Investment Wedge}

The investment Euler equation is where TFP uncertainty most naturally appears
as a second-order risk correction. Define
\begin{equation}
    m^k_{t+1}
    \equiv
    \log\beta
    +
    \log U_C(C_{t+1},L_{t+1})
    -
    \log U_C(C_t,L_t)
    +
    \log R^k_{t+1}.
    \label{eq:capital_euler_log_return}
\end{equation}
Without an investment wedge, the nonlinear Euler equation is
\begin{equation}
    \E_t[\exp(m^k_{t+1})]=1.
    \label{eq:capital_euler_no_wedge}
\end{equation}
A second-order log-normal approximation gives
\begin{equation}
    0
    \approx
    \E_t m^k_{t+1}
    +
    \frac{1}{2}\Var_t(m^k_{t+1}).
    \label{eq:capital_euler_second_order}
\end{equation}
Thus, if one writes a first-order accounting Euler equation with a log
investment wedge \(\xi^x_t\),
\begin{equation}
    0
    =
    \E_t m^k_{t+1}
    -
    \xi^x_t,
    \label{eq:first_order_investment_wedge}
\end{equation}
then the second-order no-friction model is represented by
\begin{equation}
    \xi^x_t
    =
    -\frac{1}{2}\Var_t(m^k_{t+1}).
    \label{eq:investment_uncertainty_wedge}
\end{equation}
The sign depends on the exact convention used for \(\tau_{x,t}\). The substantive
point is invariant: uncertainty in future marginal utility and future capital
returns creates a risk term in the capital Euler equation. In canonical BCA
language, this risk term is closest to the investment wedge.

\begin{remark}
This result does not mean that TFP uncertainty is literally a financial
friction. It means that, when the nonlinear risk term is forced into a
first-order BCA accounting system, the Euler-equation residual is represented as
an investment wedge.
\end{remark}
